The analysis converted MLCW deformation, cGNSS displacement, and hydraulic head records into monthly variables. Separate Bayesian ridge regression models then related each section's monthly deformation increments to hydraulic head changes, vertical surface displacement, and seasonal variation. The estimates were evaluated under three monitoring conditions that differed in when, or whether, new MLCW observations became available. These conditions included temporary delays in finalized monthly records, reduced measurement schedules with one cumulative observation every six or twelve months, and the absence of subsequent MLCW measurements after an initial calibration period.

\subsection{Preparation of model inputs}
\label{subsec:predictor_preparation}

Model-input preparation began by aligning the MLCW and cGNSS displacement records, which were collected on different dates. A deformation time series model evaluated these records at common monthly dates and produced monthly increments. The MLCW increments formed the response, while cGNSS increments, section-specific hydraulic head changes, and seasonal terms formed the predictors.

\subsubsection{Deformation time series model}
\label{subsec:deformation_model}

For each MLCW and cGNSS displacement series, the deformation time series model represented displacement as the sum of a reference value, a linear component, annual and semiannual components, and offsets at documented discontinuities. For observation time $t$, the fitted displacement was
\begin{equation}
d(t)=d_0+v(t-t_0)
+\sum_{h\in\{1,2\}}
\left[a_h\sin(2\pi h t)+b_h\cos(2\pi h t)\right]
+\sum_{j=1}^{J}c_j H(t-t_j),
\label{eq:deformation_model}
\end{equation}
\noindent where $d_0$ denotes displacement at reference time $t_0$, $v$ denotes linear velocity, $a_h$ and $b_h$ describe the periodic components, and $H(t-t_j)$ represents a step at discontinuity time $t_j$ \citep{nikolaidis_observation_2002, bevis_2014}. The fitted series were evaluated at common monthly epochs to produce aligned cumulative displacement records. For depth section $s$, the deformation increment over the month ending at epoch $k$ was

\begin{equation}
\Delta d_{s,k}
=
d_s(t_k)-d_s(t_{k-1}),
\label{eq:monthly_deformation_increment}
\end{equation}

\noindent where $t_k$ and $t_{k-1}$ denote consecutive monthly epochs. Corresponding differences in the cGNSS series yielded monthly vertical surface displacement increments.

\subsubsection{Assembly of monthly model inputs}

Each observation in the regression dataset represented one depth section during one month, with the observed deformation increment $\Delta d_{s,k}$ serving as the response. The predictors described hydraulic conditions, vertical surface displacement, and seasonal variation.

Hydraulic head observations from the Tuku wells were used directly for the depth sections represented by local records. For the remaining sections, hydraulic head at Tuku was estimated from the regional monitoring network using ordinary kriging \citep{theodossiou2006}. Current and lagged hydraulic head changes from the target section and the other five sections were then included as candidate predictors to represent hydraulic conditions throughout the monitored profile. Concurrent and lagged cGNSS increments described vertical surface displacement, while annual, semiannual, and dry-season indicator terms represented recurring seasonal variation, including a dry-season interaction with hydraulic head change. Each section model used 38 input variables derived from hydraulic head, vertical surface displacement, and seasonal variation.

A separate calibration dataset was prepared for each depth section, and an independent regression model was fitted to each dataset. Before each six-month evaluation block, predictor centering and scaling were calculated only from the observations available for calibration and were subsequently applied to the evaluation block. Previous MLCW deformation increments were not used as predictors. \Cref{tab:predictor_summary} summarizes the information represented by each predictor group, whereas \Cref{app:predictors} provides the complete predictor definitions, data sources, temporal transformations, and lag periods.

\begin{table}[htbp]
  \centering
  \small
  \renewcommand{\arraystretch}{1.15}
  \caption{Predictor groups used to estimate monthly deformation increments at Tuku.}
  \label{tab:predictor_summary}
  \begin{tabularx}{\textwidth}{@{}p{0.21\textwidth}p{0.27\textwidth}X@{}}
    \toprule
    \textbf{Predictor group} & \textbf{Information represented} & \textbf{Role in the model} \\
    \midrule
    cGNSS displacement & Current and lagged vertical surface displacement increments & Described the integrated vertical response at Tuku \\
    Target-section hydraulic head & Current and lagged hydraulic head changes for the target section & Described hydraulic conditions associated with the section being estimated \\
    Other-section hydraulic head & Current and lagged hydraulic head changes for the other five depth sections & Included as candidate predictors representing system-wide hydraulic conditions \\
    Seasonal terms & Annual and semiannual variation & Represented recurring seasonal patterns \\
    \bottomrule
  \end{tabularx}
\end{table}

\subsection{Bayesian ridge regression}
\label{subsec:bayesian_ridge}

Bayesian ridge regression was used as a probabilistic linear regression model to estimate monthly deformation increments independently for each depth section. The predictors included current and lagged changes in hydraulic head and vertical surface displacement, together with seasonal terms. Several of these predictors described related temporal variations and therefore contained overlapping information. Under these conditions, ordinary least squares may produce regression coefficients that change markedly in response to small changes in the calibration data. Ridge regression limits this sensitivity by reducing the magnitude of coefficients that are weakly supported by the observations \citep{hoerl_1970_ridge}.

The Bayesian formulation applies the same coefficient shrinkage but also represents the remaining uncertainty in the fitted relation \citep{mackay_bayesian_1992}. For $n$ calibration observations and $p$ standardized predictors, the regression model was

\begin{equation}
\Delta\boldsymbol{d}_s
=
\beta_{0,s}\boldsymbol{1}
+
\boldsymbol{X}_s\boldsymbol{\beta}_s
+
\boldsymbol{\varepsilon}_s,
\label{eq:brr_regression}
\end{equation}

\noindent with residual errors described by

\begin{equation}
\boldsymbol{\varepsilon}_s
\sim
\mathcal{N}
\left(
\boldsymbol{0},
\alpha_s^{-1}\boldsymbol{I}
\right),
\label{eq:brr_likelihood}
\end{equation}

\noindent Here, $\Delta\boldsymbol{d}_s$ contains the observed monthly deformation increments for section $s$, $\beta_{0,s}$ is the intercept, $\boldsymbol{X}_s$ contains the standardized predictors, and $\boldsymbol{\beta}_s$ contains their regression coefficients. The residual term $\boldsymbol{\varepsilon}_s$ represents variation in monthly deformation that was not explained by the predictors. The parameter $\alpha_s$ describes the precision of this residual variation. A larger value of $\alpha_s$ corresponds to less unexplained variation around the fitted relation.

Coefficient shrinkage was introduced by assigning a zero-centered Gaussian distribution to the regression coefficients,

\begin{equation}
\boldsymbol{\beta}_s
\sim
\mathcal{N}
\left(
\boldsymbol{0},
\lambda_s^{-1}\boldsymbol{I}
\right),
\label{eq:brr_prior}
\end{equation}

\noindent where $\lambda_s$ controls the strength of the shrinkage for section $s$. A larger value of $\lambda_s$ pulls weakly supported coefficients more strongly toward zero, whereas predictors that are consistently related to monthly deformation have larger coefficients. The values of $\alpha_s$ and $\lambda_s$ were estimated from the calibration data by maximizing the marginal likelihood \citep{mackay_bayesian_1992}. In practical terms, this procedure allowed the observed record to determine both the unexplained variation and the degree of coefficient shrinkage.

Conventional ridge regression also reduces unstable coefficients, but it generally provides one fitted set of coefficients for a selected penalty. Bayesian ridge regression additionally describes a range of plausible coefficient values after the calibration data have been considered. This posterior distribution provides the information needed to quantify uncertainty in subsequent monthly estimates.

A separate model was fitted for each depth section. Thus, the regression coefficients for a section were estimated only from its own MLCW deformation record, although its predictors could include hydraulic head changes from other sections of the monitored profile. The fitted relations were interpreted as statistical associations within the Tuku monitoring record and not as a substitute for a coupled groundwater flow and aquifer-system compaction model.

\subsection{Predictive uncertainty quantification}
\label{subsec:predictive_uncertainty}

The Bayesian formulation provided both a monthly deformation estimate and a measure of uncertainty around that estimate. This uncertainty was represented by the posterior predictive distribution, which combines variation not explained by the regression model with uncertainty in the fitted regression coefficients \citep{mackay_bayesian_1992}. For a new monthly epoch with standardized predictors $\boldsymbol{x}_{s,*}$ and calibration data $\mathcal{D}_s$, the posterior predictive distribution for section $s$ was

\begin{equation}
p
\left(
\Delta d_{s,*}
\mid
\boldsymbol{x}_{s,*},
\mathcal{D}_s
\right)
=
\mathcal{N}
\left(
\mu_{s,*},
\sigma_{s,*}^{2}
\right).
\label{eq:posterior_predictive}
\end{equation}

Its predictive mean was

\begin{equation}
\mu_{s,*}
=
\widehat{\beta}_{0,s}
+
\boldsymbol{x}_{s,*}^{\mathsf{T}}
\overline{\boldsymbol{\beta}}_s,
\label{eq:predictive_mean}
\end{equation}

The corresponding point estimate was

\begin{equation}
\widehat{\Delta d}_{s,*}
=
\mu_{s,*}.
\label{eq:monthly_deformation_estimate}
\end{equation}

The predictive variance was

\begin{equation}
\sigma_{s,*}^{2}
=
\alpha_s^{-1}
+
\boldsymbol{x}_{s,*}^{\mathsf{T}}
\boldsymbol{\Sigma}_{\beta,s}
\boldsymbol{x}_{s,*}.
\label{eq:predictive_variance}
\end{equation}

\noindent Here, $\widehat{\beta}_{0,s}$ is the fitted intercept, $\overline{\boldsymbol{\beta}}_s$ is the posterior mean of the regression coefficients, and $\boldsymbol{\Sigma}_{\beta,s}$ is their posterior covariance matrix. The first term in the predictive variance represents residual variation that was not explained by the model. The second term represents uncertainty in the fitted coefficients for the corresponding predictor values. The predictive interval therefore became wider when the residual variation was larger or when the coefficients were less strongly constrained by the calibration data.

The nominal 90\% posterior predictive interval was calculated as

\begin{equation}
\left[
\mu_{s,*}-1.645\sigma_{s,*},
\mu_{s,*}+1.645\sigma_{s,*}
\right],
\label{eq:prediction_interval}
\end{equation}

This interval described the range of monthly deformation increments considered plausible under the fitted model. Because it was calculated directly from the posterior predictive distribution, an interval was available for every monthly estimate without requiring an archive of previous prediction errors.

Interval performance was evaluated using empirical coverage and mean interval width. Empirical coverage measured the proportion of observed monthly increments enclosed by the intervals. Mean interval width described how broadly the plausible values were distributed around the monthly estimates. Both measures were examined because a wider interval may enclose more observations while providing less precise information \citep{singh_uncertainty_2024}. The nominal 90\% level provided a reference for this evaluation, but it did not guarantee that exactly 90\% of the observations would be enclosed. Temporal dependence and relations not represented by the regression model could cause the empirical coverage to differ from the nominal level.

\subsection{Experimental design}
\label{subsec:experimental_design}

All three evaluation designs preserved the temporal order of the monitoring records. For each period being estimated, the model used only MLCW observations that would have been available before that period, while monthly GWL and cGNSS observations remained available during the estimation period. This design prevented later MLCW observations from influencing earlier estimates. The three evaluations differed in when, or whether, new MLCW observations became available for updating the model.

\subsubsection{Monthly estimation with delayed MLCW records}
\label{subsec:delayed_evaluation}

Temporary gaps in the availability of finalized monthly MLCW records were examined using a recurring estimation and update cycle. The model was first calibrated using three years of continuous monthly observations. During each subsequent six-month period, the available GWL and cGNSS observations were used to estimate monthly deformation increments while the corresponding MLCW observations remained unavailable. The six-month interval provided a consistent basis for comparison and did not represent a fixed schedule for data provision.

At the end of each estimation period, the corresponding monthly MLCW observations became available and were compared with the estimates. These observations were then added to the calibration record, after which the model was recalibrated for the next cycle. Repeating this sequence gradually expanded the calibration record while preserving temporal order because observations from a later period could not influence earlier estimates. One complete cycle is illustrated in \Cref{fig:delayed_observation_design}.

Point estimates were evaluated separately for each depth section and for all sections combined using the coefficient of determination ($R^2$), root mean square error (RMSE), and mean absolute error (MAE). Posterior predictive intervals were evaluated using empirical coverage and mean interval width, as described in \Cref{subsec:predictive_uncertainty}. Performance was also compared across successive months within each update cycle to determine whether estimation quality changed as the MLCW observation gap lengthened.

\begin{figure}[htbp]
  \centering
  \resizebox{\linewidth}{!}{%
  \begin{tikzpicture}[
    font=\small,
    dataset/.style={anchor=east,align=right,font=\small\bfseries,text=black!75},
    phase/.style={font=\small\bfseries,text=black!75,align=center},
    updatephase/.style={font=\small\bfseries,text=orange!75!black,align=center},
    available/.style={circle,draw=blue!60!black,fill=blue!60!black,minimum size=6.3pt,inner sep=0pt},
    pending/.style={circle,draw=blue!60!black,fill=white,line width=0.9pt,minimum size=7.2pt,inner sep=0pt},
    newrecord/.style={circle,draw=orange!80!black,fill=orange!80!black,minimum size=6.6pt,inner sep=0pt},
    historyline/.style={draw=blue!55!black,line width=1.1pt},
    pendingline/.style={draw=blue!55!black,densely dashed,line width=0.9pt},
    divider/.style={draw=black!35,densely dashed,line width=0.7pt},
    flow/.style={-{Stealth[length=2.5mm]},draw=black!65,line width=0.8pt},
    orangeflow/.style={-{Stealth[length=2.5mm]},draw=orange!75!black,line width=0.9pt},
    orangeline/.style={draw=orange!80!black,densely dashed,line width=1.1pt},
    curvedflow/.style={-{Stealth[length=2.8mm]},draw=orange!80!black,line width=0.75pt},
    fadedarrow/.style={
      single arrow,
      draw=none,
      fill={rgb,255:red,34;green,139;blue,34},
      fill opacity=0.2,
      shape border rotate=90,
      minimum height=1.5cm,
      minimum width=1.0cm,
      single arrow head extend=0.8mm,
      inner sep=1pt
    },
    action/.style={draw=black!65,fill=green!8,rounded corners=2pt,align=center,minimum height=0.85cm,text width=3.05cm,font=\small\bfseries}
  ]
    \node[phase] at (2.00,3.70) {Initial historical records};
    \node[phase] at (6.00,3.70) {Estimation period};
    \node[updatephase] at (10.00,3.70) {New records};

    \node[dataset] at (-0.45,2.55) {Deformation by\\depth interval};
    \node[dataset] at (-0.45,1.55) {Hydraulic head};
    \node[dataset] at (-0.45,0.55) {Vertical surface\\displacement};

    \node[fadedarrow] at (6.00,1.55) {};

    \foreach \y in {1.55,0.55} {
      \draw[historyline] (0.40,\y) -- (7.40,\y);
    }
    \draw[historyline] (0.40,2.55) -- (4.25,2.55);
    \draw[pendingline] (4.25,2.55) -- (7.40,2.55);

    \foreach \y in {2.55,1.55,0.55} {
      \foreach \x in {0.40,1.10,1.80,3.20,3.90} {
        \node[available] at (\x,\y) {};
      }
      \node[fill=white,inner xsep=2pt,text=blue!60!black] at (2.50,\y) {$\cdots$};
    }

    \foreach \y in {1.55,0.55} {
      \foreach \x in {4.60,5.30,6.00,6.70,7.40} {
        \node[available] at (\x,\y) {};
      }
    }

    \foreach \x in {4.60,5.30,6.00,6.70,7.40} {
      \node[pending] at (\x,2.55) {};
    }

    \draw[orangeline] (8.60,2.55) -- (11.40,2.55);
    \foreach \x in {8.60,9.30,10.00,10.70,11.40} {
      \node[newrecord] at (\x,2.55) {};
    }

    \draw[curvedflow] (10.00,2.75) to[out=120,in=30,looseness=0.7] (6.00,2.75);

    \draw[divider] (4.25,-0.2) -- (4.25,4.15);
    \draw[divider] (7.85,-0.2) -- (7.85,4.15);

    \node[action] (calibrate) at (2.00,-0.75) {Calibrate\\model};
    \node[action] (estimate) at (6.00,-0.75) {Estimate\\$\widehat{\Delta d}_{s,k}$};
    \node[action] (update) at (10.00,-0.75) {Evaluate +\\recalibrate};
    \draw[flow] (calibrate.east) -- (estimate.west);
    \draw[flow] (estimate.east) -- (update.west);
    \draw[orangeflow] (7.4,2.3) to[out=-60,in=90] (update.north);
    \draw[flow,rounded corners=5pt]
      (update.south) -- ++(0,-0.68) -| (calibrate.south)
      node[pos=0.48,below,align=center,font=\small]
      {Expanded historical records\\Next cycle};
  \end{tikzpicture}%
  }
  \caption{Estimation and model updating during the temporary unavailability of finalized monthly depth-interval deformation records. Hydraulic head and vertical surface displacement remained available during this period. Filled blue markers denote available observations, open blue markers denote deformation records not yet available, and orange markers denote the records subsequently provided. The symbol $\widehat{\Delta d}_{s,k}$ denotes the estimated monthly deformation increment for section $s$. Dots indicate omitted records; marker counts and spacing are illustrative.}
  \label{fig:delayed_observation_design}
\end{figure}

\subsubsection{Sensitivity to less frequent MLCW field measurements}
\label{subsec:sparse_measurement_sensitivity}

Reduced MLCW measurement frequency was examined under hypothetical schedules in which MLCW observations became available only at the end of each measurement interval. Monthly GWL and cGNSS observations remained available throughout the interval and were used to estimate monthly deformation increments. The intervening monthly MLCW values were excluded from estimation and model updating, and one cumulative MLCW deformation observation became available at the interval endpoint.

The evaluated scenarios combined measurement intervals of 6 or 12 months with initial calibration periods containing 36, 60, or 96 consecutive monthly MLCW observations. All three calibration periods ended on the same calendar date, so the 36- and 60-month periods began later than the 96-month period rather than starting from the same date and ending earlier. This shared calibration endpoint let the three initial-history lengths be compared over an identical subsequent record rather than over three different sets of calendar months. Each measurement schedule was then evaluated over an identical subsequent period, ending at a shared field-measurement date, so that the three initial-history lengths underwent the same number of model updates at the same calendar dates within a given measurement schedule. Because a 6-month schedule and a 12-month schedule complete different numbers of measurement intervals over the same period, the two schedules were compared using their own respective update counts rather than a common count; only the initial-history-length comparison required matching update counts and dates. These durations defined the scenarios examined in this study rather than a requirement of the method. Regardless of the selected duration, each measurement interval produced one cumulative deformation observation.

To represent this cumulative observation mathematically, let $s$ denote a depth section, $k$ denote a monthly epoch, $I$ denote a measurement interval, and $H_I$ denote the number of months in that interval. The observed monthly MLCW deformation increment is denoted by $\Delta d_{s,k}$, as defined in \Cref{eq:monthly_deformation_increment}. The cumulative deformation observed over interval $I$ was

\begin{equation}
\Delta d_{s,I}
=
\sum_{k\in I}\Delta d_{s,k}.
\label{eq:cumulative_deformation_observation}
\end{equation}

The individual monthly values within $I$ remained hidden from model fitting and updating and were used only for retrospective evaluation. Consequently, only the cumulative observation $\Delta d_{s,I}$ entered the next model update. Representing this observation within the monthly regression required the relation in \Cref{eq:brr_regression,eq:brr_likelihood} to be accumulated over the same interval. This summation gave

\begin{equation}
\Delta d_{s,I}
=
H_I\beta_{0,s}
+
\left(
\sum_{k\in I}\boldsymbol{x}_{s,k}
\right)^{\mathsf{T}}
\boldsymbol{\beta}_{s}
+
\varepsilon_{s,I},
\label{eq:cumulative_deformation_model}
\end{equation}

\noindent where $\boldsymbol{x}_{s,k}$ contains the standardized predictors for section $s$ at monthly epoch $k$, and $\varepsilon_{s,I}$ is the sum of the monthly residuals within the interval. This cumulative equation assumed that the monthly relation between the predictors and deformation remained unchanged between successive model updates. It also preserved the residual model in \Cref{eq:brr_likelihood}, under which the monthly residuals were independent, normally distributed around zero, and described by a constant variance within each depth section. Their variances therefore accumulated over the measurement interval, and the cumulative residual followed

\begin{equation}
\varepsilon_{s,I}
\sim
\mathcal{N}
\left(
0,
H_I\alpha_s^{-1}
\right).
\label{eq:cumulative_residual_distribution}
\end{equation}

The resulting variance increased in proportion to $H_I$. To fit the cumulative observation together with the original monthly observations, both observation types were placed on the same residual variance scale. The monthly predictors were first centered and scaled using only the initial monthly calibration record. The cumulative regression equation was then divided by $\sqrt{H_I}$,

\begin{equation}
\frac{\Delta d_{s,I}}{\sqrt{H_I}}
=
\sqrt{H_I}\,\beta_{0,s}
+
\left(
\frac{1}{\sqrt{H_I}}
\sum_{k\in I}\boldsymbol{x}_{s,k}
\right)^{\mathsf{T}}
\boldsymbol{\beta}_{s}
+
\widetilde{\varepsilon}_{s,I},
\label{eq:scaled_cumulative_deformation_model}
\end{equation}

\noindent with

\begin{equation}
\widetilde{\varepsilon}_{s,I}
=
\frac{\varepsilon_{s,I}}{\sqrt{H_I}}
\sim
\mathcal{N}
\left(
0,
\alpha_s^{-1}
\right).
\label{eq:scaled_cumulative_residual}
\end{equation}

This scaling did not divide the cumulative observation into separate monthly measurements. Instead, it allowed the monthly and cumulative observations to constrain the same regression coefficients while preserving their different temporal support. The scaled cumulative equation could therefore enter the calibration record as one interval constraint.

The model was first calibrated from the initial monthly MLCW record and used to estimate deformation throughout the first measurement interval. At the end of that interval, the cumulative observation in \Cref{eq:cumulative_deformation_observation} was added to the calibration record as one interval constraint. The expanded record, consisting of the initial monthly observations and cumulative observations from completed intervals, was then used to recalibrate the model before the next interval. Hidden monthly MLCW values from the reduced measurement intervals never entered this recalibration. This sequence is illustrated in \Cref{fig:reduced_measurement_design}.

\begin{figure}[htbp]
  \centering
  \resizebox{\linewidth}{!}{%
  \begin{tikzpicture}[
    x=1.12cm,
    y=1cm,
    font=\small,
    dataset/.style={anchor=east,align=right,font=\small\bfseries,text=black!75},
    phase/.style={font=\small\bfseries,text=black!75,align=center},
    available/.style={circle,draw=blue!60!black,fill=blue!60!black,minimum size=6.3pt,inner sep=0pt},
    estimate/.style={circle,draw=black!55,fill=white,line width=1.0pt,minimum size=7.2pt,inner sep=0pt},
    startrecord/.style={circle,draw=blue!60!black,fill=blue!60!black,minimum size=8.5pt,inner sep=0pt},
    latestrecord/.style={circle,draw=orange!80!black,fill=orange!80!black,minimum size=9pt,inner sep=0pt},
    historyline/.style={draw=blue!55!black,line width=1.1pt},
    estimateline/.style={draw=black!50,densely dashed,line width=0.9pt},
    divider/.style={draw=black!35,densely dashed,line width=0.7pt},
    flow/.style={-{Stealth[length=2.5mm]},draw=black!65,line width=0.8pt},
    predictorflow/.style={-{Stealth[length=2.7mm]},draw=green!45!black,line width=1.15pt},
    orangeflow/.style={-{Stealth[length=2.5mm]},draw=orange!75!black,line width=0.9pt},
    toparrow/.style={{Stealth[length=3mm]}-{Stealth[length=3mm]},draw=black,line width=1.4pt},
    fadedarrow/.style={
      single arrow,
      draw=none,
      fill={rgb,255:red,34;green,139;blue,34},
      fill opacity=0.2,
      shape border rotate=90,
      minimum height=1.5cm,
      minimum width=1.0cm,
      single arrow head extend=0.8mm,
      inner sep=1pt
    },
    action/.style={draw=black!65,fill=green!8,rounded corners=2pt,align=center,minimum height=0.85cm,text width=3.05cm,font=\small\bfseries}
  ]
    \node[phase] at (2.00,4.5) {Initial historical records};
    \node[phase] at (6.00,4.5) {Estimation period};

    \node[dataset] at (-0.45,2.55) {Deformation by\\depth interval};
    \node[dataset] at (-0.45,1.55) {Hydraulic head};
    \node[dataset] at (-0.45,0.55) {Vertical surface\\displacement};

    \node[fadedarrow] at (6.00,1.55) {};

    \foreach \y in {1.55,0.55} {
      \draw[historyline] (0.40,\y) -- (7.40,\y);
    }
    \draw[historyline] (0.40,2.55) -- (3.90,2.55);
    \draw[estimateline] (3.90,2.55) -- (7.40,2.55);

    \foreach \y in {1.55,0.55} {
      \foreach \x in {0.40,1.10,1.80,3.20,3.90} {
        \node[available] at (\x,\y) {};
      }
      \node[fill=white,inner xsep=2pt,text=blue!60!black] at (2.50,\y) {$\cdots$};
    }
    \foreach \x in {0.40,1.10,1.80,3.20} {
      \node[available] at (\x,2.55) {};
    }
    \node[fill=white,inner xsep=2pt,text=blue!60!black] at (2.50,2.55) {$\cdots$};
    \node[startrecord] (start) at (3.90,2.55) {};

    \foreach \y in {1.55,0.55} {
      \foreach \x in {4.60,5.30,6.00,6.70,7.40} {
        \node[available] at (\x,\y) {};
      }
    }
    \foreach \x in {4.60,5.30,6.00,6.70} {
      \node[estimate] at (\x,2.55) {};
    }

    \node[latestrecord] (finish) at (7.40,2.55) {};

    \draw[toparrow]
      (start.north) -- ++(0,0.85) -| (finish.north)
      node[pos=0.25,above=-1.5pt,font=\small,text=black] (intervaldeformation)
      {$\Delta d_{s,I}$};

    \draw[divider] (4.25,-0.2) -- (4.25,4.35);
    \draw[divider] (7.85,-0.2) -- (7.85,4.35);

    \node[action] (calibrate) at (2.00,-0.75) {Calibrate\\model};
    \node[action] (estimatebox) at (6.00,-0.75) {Estimate\\$\widehat{\Delta d}_{s,k}$};
    \node[action] (update) at (10.00,-0.75) {Evaluate +\\recalibrate};
    \draw[flow] (calibrate.east) -- (estimatebox.west);
    \draw[flow] (estimatebox.east) -- (update.west);
    \draw[orangeflow] (intervaldeformation.east) -| (update.north);

    \draw[flow,rounded corners=5pt]
      (update.south) -- ++(0,-0.68) -| (calibrate.south)
      node[pos=0.48,below,align=center,font=\small]
      {Expanded historical records\\Next cycle};
  \end{tikzpicture}%
  }
  \caption{Monthly deformation estimation under reduced measurement frequency. Filled blue markers denote available observations, while hydraulic head and vertical surface displacement remained available monthly. Open gray markers denote monthly estimates $\widehat{\Delta d}_{s,k}$, and the blue and orange endpoints define the observed cumulative deformation $\Delta d_{s,I}$ over interval $I$. Dots indicate omitted records; marker counts and spacing are illustrative.}
  \label{fig:reduced_measurement_design}
\end{figure}

The resulting estimates were evaluated at monthly and interval scales. At the monthly scale, the error compared each estimated increment with its hidden monthly MLCW value,

\begin{equation}
e_{s,k}
=
\widehat{\Delta d}_{s,k}
-
\Delta d_{s,k}.
\label{eq:sparse_monthly_error}
\end{equation}

At the interval scale, the endpoint error compared the accumulated monthly estimates with the cumulative observation,

\begin{equation}
e_{s,I}
=
\sum_{k\in I}\widehat{\Delta d}_{s,k}
-
\Delta d_{s,I}.
\label{eq:sparse_endpoint_error}
\end{equation}

Monthly errors described how well the model distributed deformation within each interval, whereas endpoint errors described agreement with the cumulative field observation. Both error types were summarized using MAE, RMSE, and mean signed error for each depth section and for all sections combined. The coefficient of determination was not calculated for the endpoint observations because each scenario contained relatively few cumulative measurements.

\subsubsection{Sensitivity without subsequent MLCW field measurements}
\label{subsec:permanent_stoppage_sensitivity}

The absence of subsequent MLCW measurements was examined using models calibrated once from initial records spanning 3, 5, or 8 years. All three calibration records ended on the same calendar date, so the 3- and 5-year records began later than the 8-year record rather than starting from the same date and ending earlier. This shared calibration endpoint let the three initial-history lengths be compared over an identical subsequent record rather than over three different sets of calendar months. After each calibration record ended, monthly GWL and cGNSS observations continued to provide predictor values, but no later MLCW observations were added. The model coefficients and predictor scaling therefore remained fixed throughout the subsequent estimation period.

The remaining historical MLCW observations allowed this hypothetical condition to be evaluated retrospectively. These observations were withheld from model fitting and were used only to compare the observed monthly deformation increments with the estimates generated after the calibration record. The monthly error $e_{s,k}$ was calculated as defined in \Cref{eq:sparse_monthly_error}.

The accumulation of monthly errors was evaluated over the first $K$ months after calibration. The cumulative signed error for section $s$ was

\begin{equation}
E_{s,K}
=
\sum_{k=1}^{K} e_{s,k},
\label{eq:stoppage_cumulative_signed_error}
\end{equation}

and its absolute magnitude was

\begin{equation}
A_{s,K}
=
\left|E_{s,K}\right|.
\label{eq:stoppage_absolute_cumulative_error}
\end{equation}

The signed value $E_{s,K}$ indicated whether the accumulated estimates were above or below the observed cumulative deformation, whereas $A_{s,K}$ described the magnitude of this difference without regard to direction.

Monthly estimation performance through horizon $K$ was summarized separately using

\begin{equation}
\operatorname{MAE}_{s,K}
=
\frac{1}{K}
\sum_{k=1}^{K}
\left|e_{s,k}\right|,
\label{eq:stoppage_running_mae}
\end{equation}

and

\begin{equation}
\operatorname{RMSE}_{s,K}
=
\sqrt{
\frac{1}{K}
\sum_{k=1}^{K}
e_{s,k}^{2}
}.
\label{eq:stoppage_running_rmse}
\end{equation}

These measures described the typical magnitude of the monthly errors from the end of calibration through horizon $K$. They therefore complemented $A_{s,K}$, which described the difference between the estimated and observed cumulative deformation at that horizon.

To account for the increasing duration of the estimation period, the absolute cumulative error was also expressed per elapsed month as

\begin{equation}
R_{s,K}
=
\frac{A_{s,K}}{K}.
\label{eq:stoppage_horizon_normalized_error}
\end{equation}

The normalized value $R_{s,K}$ was interpreted together with $A_{s,K}$. Because positive and negative monthly errors could partly cancel within $E_{s,K}$, a decrease in $R_{s,K}$ did not necessarily indicate that the absolute cumulative error had stopped growing.

The frozen Bayesian ridge model also provided the posterior predictive interval defined in \Cref{subsec:predictive_uncertainty} for each monthly estimate. These intervals described uncertainty in individual monthly deformation increments throughout the estimation period. They were not accumulated into uncertainty bands for $E_{s,K}$ or $A_{s,K}$ because the dependence among prediction errors from successive months was not represented in the cumulative calculation.

Because the three calibration records shared the same endpoint, all three left the same 80 months for retrospective evaluation, over the identical subsequent calendar period. The complete fit-once evaluation is illustrated in \Cref{fig:no_subsequent_mlcw_design}.

\begin{figure}[htbp]
  \centering
  \resizebox{\linewidth}{!}{%
  \begin{tikzpicture}[
  	x=1.12cm,
  	y=1cm,
  	font=\small,
  	dataset/.style={anchor=east,align=right,font=\small\bfseries,text=black!75},
  	phase/.style={font=\small\bfseries,text=black!75,align=center},
  	available/.style={circle,draw=blue!60!black,fill=blue!60!black,minimum size=6.3pt,inner sep=0pt},
  	estimate/.style={circle,draw=black!55,fill=white,line width=1.0pt,minimum size=7.2pt,inner sep=0pt},
  	historyline/.style={draw=blue!55!black,line width=1.1pt},
  	estimateline/.style={draw=black!50,densely dashed,line width=0.9pt},
  	divider/.style={draw=black!35,densely dashed,line width=0.7pt},
  	flow/.style={-{Stealth[length=2.5mm]},draw=black!65,line width=0.8pt},
  	fadedarrow/.style={
  		single arrow,
  		draw=none,
  		fill={rgb,255:red,34;green,139;blue,34},
  		fill opacity=0.2,
  		shape border rotate=90,
  		minimum height=1.5cm,
  		minimum width=1.0cm,
  		single arrow head extend=0.8mm,
  		inner sep=1pt
  	},
  	action/.style={draw=black!65,fill=green!8,rounded corners=2pt,align=center,minimum height=0.85cm,text width=3.05cm,font=\small\bfseries}
  	]
  	\node[phase] at (2.00,3.70) {Initial historical records};
  	\node[phase] at (6.20,3.70) {Estimation period};
  	
  	\node[dataset] at (-0.45,2.55) {Deformation by\\depth interval};
  	\node[dataset] at (-0.45,1.55) {Hydraulic head};
  	\node[dataset] at (-0.45,0.55) {Vertical surface\\displacement};
  	
  	% Continuous timelines
  	\foreach \y in {1.55,0.55} {
  		\draw[historyline] (0.40,\y) -- (8.00,\y);
  	}
  	\draw[historyline] (0.40,2.55) -- (3.90,2.55);
  	\draw[estimateline] (3.90,2.55) -- (8.00,2.55);
  	
  	% Panel 1 observations
  	\foreach \y in {2.55,1.55,0.55} {
  		\foreach \x in {0.40,1.10,1.80,3.20,3.90} {
  			\node[available] at (\x,\y) {};
  		}
  		\node[fill=white,inner xsep=2pt,text=blue!60!black] at (2.50,\y) {$\cdots$};
  	}
  	
  	% Panel 2 observations & dots
  	\foreach \y in {1.55,0.55} {
  		\foreach \x in {4.60,5.30,7.30,8.00} {
  			\node[available] at (\x,\y) {};
  		}
  		\node[fill=white,inner xsep=2pt,text=blue!60!black] at (6.30,\y) {$\cdots$};
  	}
  	\foreach \x in {4.60,5.30,7.30,8.00} {
  		\node[estimate] at (\x,2.55) {};
  	}
  	\node[fill=white,inner xsep=2pt,text=black!55] at (6.30,2.55) {$\cdots$};
  	
  	% CHUYỂN MŨI TÊN XANH XUỐNG ĐÂY VÀ ĐỔI SANG X=5.30
  	\node[fadedarrow] at (6.30,1.55) {};
  	
  	\draw[divider] (4.25,-0.2) -- (4.25,4.15);
  	
  	\node[action] (calibrate) at (2.00,-0.75) {Calibrate\\model};
  	\node[action] (estimatebox) at (6.20,-0.75) {Estimate\\$\widehat{\Delta d}_{s,k}$};
  	\draw[flow] (calibrate.east) -- (estimatebox.west);
  	\node[font=\small,text=black!70] at (6.20,-1.55) {No recalibration};
  \end{tikzpicture}
  }
  \caption{Fit-once monthly deformation estimation under the hypothetical absence of subsequent depth-interval deformation measurements. Filled blue markers denote observations used for model calibration. Hydraulic head and vertical surface displacement remained available, while open gray markers denote monthly estimates $\widehat{\Delta d}_{s,k}$. Model coefficients and predictor scaling remained fixed throughout the estimation period. Dots indicate omitted records; marker counts and spacing are illustrative.}
  \label{fig:no_subsequent_mlcw_design}
\end{figure}

\FloatBarrier
